% ch01_introduction.tex — Introduction et Motivations
\chapter{Introduction et Motivations}\label{ch:introduction}

Consid\'erez un barman qui m\'elange un cocktail dans un verre. Apr\`es avoir
vigoureusement agit\'e le liquide, il est certain qu'au moins un point du
liquide se retrouve exactement \`a sa position initiale. Ce r\'esultat,
connu sous le nom de th\'eor\`eme de Brouwer, illustre une id\'ee simple mais
profonde~: sous certaines conditions, toute transformation poss\`ede au
moins un \emph{point fixe}, un point qui ne bouge pas. Cette id\'ee traverse
les math\'ematiques~: elle appara\^it en analyse, en topologie, en
alg\`ebre, en \'economie (l'\'equilibre de Nash), en informatique (la
s\'emantique des programmes), et en th\'eorie des jeux. Ce cours explore
syst\'ematiquement les th\'eor\`emes de points fixes, leurs d\'emonstrations,
et leurs applications.

% ------------------------------------------------------------
\section{Qu'est-ce qu'un point fixe~?}
% ------------------------------------------------------------

La d\'efinition est d'une simplicit\'e d\'esarmante --- mais ne vous y fiez pas.

\begin{definition}[Point fixe]
Soit $X$ un ensemble et $T : X \to X$ une application. Un \'el\'ement $x^* \in X$
est un \emph{point fixe} de $T$ si $T(x^*) = x^*$.
L'ensemble des points fixes de $T$ est not\'e
\[
    \mathrm{Fix}(T) = \{ x \in X : T(x) = x \}.
\]
\end{definition}

Cette d\'efinition, d'apparence \'el\'ementaire, conduit \`a des questions
remarquablement profondes d\`es que l'on enrichit $X$ d'une structure
suppl\'ementaire (m\'etrique, topologie, ordre) et que l'on impose des
conditions sur $T$.

\begin{remarque}
L'\'equation $T(x) = x$ peut \^etre reformul\'ee comme $F(x) = 0$ en posant
$F = T - \mathrm{Id}$. R\'eciproquement, toute \'equation $F(x) = 0$ peut se
ramener \`a un probl\`eme de point fixe en posant $T(x) = x - F(x)$ (ou
$T(x) = x + \alpha F(x)$ pour un param\`etre $\alpha$ bien choisi). Cette
dualit\'e est \`a la base de nombreuses m\'ethodes num\'eriques.
\end{remarque}

\section{Contexte historique}

La th\'eorie des points fixes a une histoire riche et fascinante.

\subsection{Les origines~: Picard et Banach}

La m\'ethode des approximations successives remonte aux travaux d'\'Emile Picard
(1890) sur les \'equations diff\'erentielles ordinaires. Picard a montr\'e que
l'\'equation
\[
    y'(t) = f(t, y(t)), \quad y(t_0) = y_0
\]
poss\`ede une solution unique sous une condition de Lipschitz sur $f$ par rapport
\`a $y$, en construisant la suite d'it\'er\'ees
\[
    y_{n+1}(t) = y_0 + \int_{t_0}^{t} f(s, y_n(s)) \, ds.
\]

C'est Stefan Banach qui, dans sa th\`ese de doctorat (1920, publi\'ee en 1922),
a formalis\'e cette id\'ee en \'enon\c{c}ant le \emph{Principe de Contraction} dans
le cadre g\'en\'eral des espaces m\'etriques complets.

\subsection{Le th\'eor\`eme de Brouwer}

En 1911, L.\,E.\,J.\ Brouwer a d\'emontr\'e qu'en dimension finie, toute
application continue d'un convexe compact dans lui-m\^eme poss\`ede un point fixe.
Ce r\'esultat, de nature purement topologique (aucune hypoth\`ese m\'etrique),
repose sur des arguments de degr\'e topologique et de topologie alg\'ebrique.

\subsection{Extensions en dimension infinie}

Juliusz Schauder (1930) a g\'en\'eralis\'e le th\'eor\`eme de Brouwer aux espaces de
Banach de dimension infinie, en rempla\c{c}ant la compacit\'e de l'ensemble par
la compacit\'e de l'application. Shizuo Kakutani (1941) a \'etendu le r\'esultat
aux correspondances multivoques (applications \`a valeurs ensemblistes).

\subsection{Approche par l'ordre}

Alfred Tarski (1955) a \'etabli un th\'eor\`eme de point fixe dans les treillis
complets, fond\'e uniquement sur la monotonie de l'application, sans aucune
hypoth\`ese m\'etrique ou topologique. Ce r\'esultat, souvent appel\'e
th\'eor\`eme de Knaster--Tarski, trouve des applications en informatique
th\'eorique et en logique.

\section{Exemples \'el\'ementaires}

\begin{exemple}[Point fixe en dimension 1]
Soit $f : [0, 1] \to [0, 1]$ une application continue. Par le th\'eor\`eme des
valeurs interm\'ediaires, $f$ poss\`ede un point fixe.

En effet, posons $g(x) = f(x) - x$. Alors $g(0) = f(0) \geq 0$ et
$g(1) = f(1) - 1 \leq 0$. Par le th\'eor\`eme des valeurs interm\'ediaires,
il existe $x^* \in [0,1]$ tel que $g(x^*) = 0$, c'est-\`a-dire $f(x^*) = x^*$.
\end{exemple}

\begin{exemple}[Non-existence de point fixe]
L'application $T : \R \to \R$ d\'efinie par $T(x) = x + 1$ n'a pas de point fixe.
Cet exemple montre la n\'ecessit\'e de travailler dans un ensemble
\og{}suffisamment petit\fg (born\'e, compact, etc.).
\end{exemple}

\begin{exemple}[Multiplicit\'e des points fixes]
L'application identit\'e $T(x) = x$ a tous les points comme points fixes.
L'application $T(x) = x^2$ sur $\R$ a exactement deux points fixes~: $0$ et $1$.
\end{exemple}

\begin{exemple}[Rotation sans point fixe]
La rotation $R_\theta : S^1 \to S^1$ d'angle $\theta \neq 0 \pmod{2\pi}$ sur le
cercle unit\'e n'a pas de point fixe. Ceci montre que la convexit\'e est essentielle
dans le th\'eor\`eme de Brouwer~: le cercle $S^1$ est compact mais non convexe.
\end{exemple}

\section{Formulation comme \'equation fonctionnelle}

De nombreux probl\`emes math\'ematiques classiques se reformulent comme des
probl\`emes de points fixes.

\subsection{\'Equations int\'egrales}

L'\'equation int\'egrale de Fredholm de seconde esp\`ece
\[
    u(t) = f(t) + \lambda \int_a^b K(t,s) \, u(s) \, ds
\]
est de la forme $u = Tu$ o\`u l'op\'erateur $T$ est d\'efini par
\[
    (Tu)(t) = f(t) + \lambda \int_a^b K(t,s) \, u(s) \, ds.
\]

\subsection{\'Equations diff\'erentielles}

Le probl\`eme de Cauchy $y' = f(t,y)$, $y(t_0) = y_0$ est \'equivalent, sous
des hypoth\`eses de r\'egularit\'e, au probl\`eme de point fixe $y = Ty$ o\`u
\[
    (Ty)(t) = y_0 + \int_{t_0}^t f(s, y(s)) \, ds.
\]

\subsection{\'Equations aux d\'eriv\'ees partielles}

Les probl\`emes aux limites pour les EDP elliptiques se reformulent souvent
comme des probl\`emes de points fixes dans des espaces de Sobolev. Par exemple,
le probl\`eme
\[
    -\Delta u = g(u) \text{ dans } \Omega, \quad u = 0 \text{ sur } \partial\Omega
\]
s'\'ecrit $u = (-\Delta)^{-1} g(u) = T(u)$.

\subsection{\'Equilibres de Nash}

Un \'equilibre de Nash dans un jeu \`a $n$ joueurs peut \^etre vu comme un point
fixe d'une correspondance de meilleure r\'eponse. Si $S_i$ est l'ensemble des
strat\'egies du joueur $i$ et $BR_i(s_{-i})$ sa meilleure r\'eponse aux strat\'egies
des autres joueurs, alors un \'equilibre de Nash est un profil $s^* = (s_1^*, \ldots, s_n^*)$
tel que $s_i^* \in BR_i(s_{-i}^*)$ pour tout $i$.

\section{Trois familles de th\'eor\`emes}

Les th\'eor\`emes de points fixes se r\'epartissent en trois grandes familles,
selon la structure math\'ematique utilis\'ee.

\begin{intuition}[title=\textbf{Les trois approches}]
\begin{enumerate}
    \item \textbf{Approche m\'etrique}~: on suppose que $T$ contracte les distances.
          Le point fixe est unique et obtenu par it\'eration. \\
          Prototype~: th\'eor\`eme de Banach.
    \item \textbf{Approche topologique}~: on suppose la continuit\'e de $T$ et la
          compacit\'e du domaine. L'existence est garantie, mais pas l'unicit\'e
          ni la constructibilit\'e. \\
          Prototype~: th\'eor\`eme de Brouwer.
    \item \textbf{Approche par l'ordre}~: on suppose la monotonie de $T$ dans un
          treillis complet. L'existence est garantie, et l'ensemble des points
          fixes forme lui-m\^eme un treillis complet. \\
          Prototype~: th\'eor\`eme de Tarski.
\end{enumerate}
\end{intuition}

\section{Rappels et pr\'eliminaires}

\subsection{Espaces m\'etriques}

\begin{definition}[Espace m\'etrique]
Un \emph{espace m\'etrique} est un couple $(X, d)$ o\`u $X$ est un ensemble et
$d : X \times X \to [0, +\infty)$ est une application v\'erifiant pour tous
$x, y, z \in X$~:
\begin{enumerate}
    \item $d(x, y) = 0 \Leftrightarrow x = y$ \quad (s\'eparation),
    \item $d(x, y) = d(y, x)$ \quad (sym\'etrie),
    \item $d(x, z) \leq d(x, y) + d(y, z)$ \quad (in\'egalit\'e triangulaire).
\end{enumerate}
\end{definition}

\begin{definition}[Compl\'etude]
Un espace m\'etrique $(X, d)$ est \emph{complet} si toute suite de Cauchy dans
$X$ converge. Rappelons qu'une suite $(x_n)$ est de Cauchy si pour tout
$\varepsilon > 0$, il existe $N \in \N$ tel que $d(x_m, x_n) < \varepsilon$
pour tous $m, n \geq N$.
\end{definition}

\begin{definition}[Application lipschitzienne]
Soient $(X, d_X)$ et $(Y, d_Y)$ deux espaces m\'etriques. Une application
$T : X \to Y$ est \emph{lipschitzienne} de constante $L \geq 0$ si
\[
    d_Y(T(x), T(y)) \leq L \, d_X(x, y) \quad \forall x, y \in X.
\]
Si $L < 1$, on dit que $T$ est une \emph{contraction}. Si $L = 1$, on dit
que $T$ est \emph{non expansive}.
\end{definition}

\subsection{Espaces de Banach}

\begin{definition}[Espace de Banach]
Un \emph{espace de Banach} est un espace vectoriel norm\'e $(E, \norm{\cdot})$
complet pour la distance induite $d(x,y) = \norm{x - y}$.
\end{definition}

\begin{definition}[Compacit\'e]
Un sous-ensemble $K$ d'un espace topologique est \emph{compact} si de tout
recouvrement ouvert de $K$ on peut extraire un sous-recouvrement fini.
Dans un espace m\'etrique, ceci \'equivaut \`a la compacit\'e s\'equentielle~:
toute suite dans $K$ poss\`ede une sous-suite convergente dans $K$.
\end{definition}

\begin{remarque}
En dimension infinie, les boules ferm\'ees born\'ees ne sont jamais compactes
(th\'eor\`eme de Riesz). C'est pourquoi les th\'eor\`emes topologiques en
dimension infinie (Schauder) n\'ecessitent des hypoth\`eses de compacit\'e sur
l'op\'erateur plut\^ot que sur le domaine.
\end{remarque}

\subsection{Convexit\'e}

\begin{definition}[Ensemble convexe]
Un sous-ensemble $C$ d'un espace vectoriel $E$ est \emph{convexe} si pour tous
$x, y \in C$ et tout $t \in [0,1]$, on a $(1-t)x + ty \in C$.
\end{definition}

\begin{definition}[Enveloppe convexe]
L'\emph{enveloppe convexe} d'un ensemble $A \subset E$ est le plus petit
convexe contenant $A$~:
\[
    \mathrm{co}(A) = \left\{ \sum_{i=1}^n \lambda_i x_i : n \in \N^*,
    \, x_i \in A, \, \lambda_i \geq 0, \, \sum_{i=1}^n \lambda_i = 1 \right\}.
\]
\end{definition}

\section{It\'eration de Picard}

La m\'ethode la plus \'el\'ementaire pour trouver un point fixe est l'it\'eration
de Picard (ou m\'ethode des approximations successives).

\begin{algorithme}[title=\textbf{It\'eration de Picard}]
\textbf{Entr\'ee}~: Application $T : X \to X$, point initial $x_0 \in X$. \\
\textbf{Proc\'edure}~: D\'efinir la suite $(x_n)$ par $x_{n+1} = T(x_n)$ pour $n \geq 0$. \\
\textbf{Sortie}~: Si $(x_n)$ converge vers $x^*$ et si $T$ est continue, alors $x^*$ est un point fixe de $T$.
\end{algorithme}

\begin{remarque}
La convergence de l'it\'eration de Picard n'est pas garantie en g\'en\'eral.
Le th\'eor\`eme de Banach (Chapitre~2) fournit des conditions suffisantes.
\end{remarque}

\section{Le th\'eor\`eme du point fixe en dimension 1}

\begin{theoreme}[Th\'eor\`eme du point fixe en dimension 1]
Soit $f : [a,b] \to [a,b]$ une application continue. Alors $f$ poss\`ede au
moins un point fixe dans $[a,b]$.
\end{theoreme}

\begin{proof}
Consid\'erons la fonction $g : [a,b] \to \R$ d\'efinie par $g(x) = f(x) - x$.
Puisque $f([a,b]) \subset [a,b]$, nous avons $g(a) = f(a) - a \geq 0$ et
$g(b) = f(b) - b \leq 0$. Par le th\'eor\`eme des valeurs interm\'ediaires,
il existe $x^* \in [a,b]$ tel que $g(x^*) = 0$, c'est-\`a-dire $f(x^*) = x^*$.
\end{proof}

\begin{attention}[title=\textbf{Hypoth\`eses n\'ecessaires}]
Les deux hypoth\`eses sont essentielles~:
\begin{itemize}
    \item \textbf{Continuit\'e}~: l'application $f : [0,1] \to [0,1]$ d\'efinie par
          $f(x) = 0$ si $x > 1/2$ et $f(x) = 1$ si $x \leq 1/2$ n'a pas de point fixe.
    \item \textbf{Stabilit\'e} ($f([a,b]) \subset [a,b]$)~: l'application
          $f(x) = x + 1$ sur $[0,1]$ n'a pas de point fixe car $f([0,1]) = [1,2] \not\subset [0,1]$.
\end{itemize}
\end{attention}

\section{Degr\'e topologique~: aper\c{c}u}

Le degr\'e topologique est un invariant alg\'ebrique qui \og{}compte\fg les
solutions d'une \'equation $F(x) = p$ avec multiplicit\'e et signe. Pour une
application $F : \overline{\Omega} \to \R^n$ de classe $C^1$, avec $\Omega$ un
ouvert born\'e de $\R^n$ et $p \notin F(\partial\Omega)$, le degr\'e de Brouwer
est d\'efini par
\[
    \deg(F, \Omega, p) = \sum_{x \in F^{-1}(p)} \mathrm{sgn}\,\det(DF(x)).
\]

\begin{theoreme}[Existence par degr\'e non nul]
Si $\deg(F, \Omega, p) \neq 0$, alors l'\'equation $F(x) = p$ poss\`ede au moins
une solution dans $\Omega$.
\end{theoreme}

Ce lien entre topologie alg\'ebrique et existence de solutions sera d\'evelopp\'e
en d\'etail au Chapitre~5.

\section{Plan du cours}

\begin{formules}[title=\textbf{Architecture du cours}]
\begin{center}
\renewcommand{\arraystretch}{1.3}
\begin{tabular}{|c|l|l|}
\hline
\textbf{Ch.} & \textbf{Titre} & \textbf{R\'esultat principal} \\
\hline
1 & Introduction & Motivations, rappels \\
2 & Banach & Principe de contraction \\
3 & Extensions m\'etriques & Edelstein, Kannan, \'Ciri\'c, Caristi \\
4 & Espaces g\'en\'eralis\'es & $b$-m\'etriques, $G$-m\'etriques \\
5 & Brouwer & Point fixe topologique (dim.\ finie) \\
6 & Schauder & Point fixe topologique (dim.\ infinie) \\
7 & Kakutani & Correspondances multivoques \\
8 & Espaces ordonn\'es & Tarski--Knaster \\
9 & Applications & EDO, EDP, jeux, \'economie \\
10 & Markov--Kakutani & Points fixes communs \\
11 & Points fixes al\'eatoires & Th\'eorie mesurable \\
\hline
\end{tabular}
\end{center}
\end{formules}

\section{Exercices}

\begin{exercice}
Soit $f : [0,1] \to [0,1]$ une application croissante (non n\'ecessairement
continue). Montrer que $f$ poss\`ede un point fixe.
\textit{Indication~: consid\'erer $x^* = \sup\{x \in [0,1] : f(x) \geq x\}$.}
\end{exercice}

\begin{exercice}
Soit $f : \R \to \R$ continue telle que $\abs{f(x) - f(y)} < \abs{x - y}$ pour
tous $x \neq y$. Montrer que $f$ a au plus un point fixe. Donner un exemple o\`u
$f$ n'a pas de point fixe.
\end{exercice}

\begin{exercice}
Soit $T : \R^n \to \R^n$ une isom\'etrie affine (c'est-\`a-dire $\norm{T(x) - T(y)} = \norm{x - y}$
pour tous $x, y$). Montrer que soit $T$ a un point fixe, soit $T$ est une
translation sans point fixe.
\end{exercice}

\begin{exercice}
Montrer que la rotation de $\R^2$ d'angle $\pi/4$ autour de l'origine a un
unique point fixe. Plus g\'en\'eralement, que peut-on dire des points fixes d'une
rotation d'angle $\theta$ en dimension $n$~?
\end{exercice}

\begin{exercice}
Soit $(X, d)$ un espace m\'etrique compact et $T : X \to X$ une isom\'etrie
(c'est-\`a-dire $d(T(x), T(y)) = d(x,y)$ pour tous $x,y$). Montrer que $T$
est surjective. En d\'eduire que si $X$ est de plus convexe dans un espace
norm\'e, alors $T$ poss\`ede un point fixe.
\end{exercice}

\begin{exercice}
Soit $A$ une matrice $n \times n$ \`a coefficients positifs ou nuls telle que la
somme des \'el\'ements de chaque colonne vaut $1$ (matrice stochastique).
Montrer que $A$ poss\`ede un vecteur propre associ\'e \`a la valeur propre $1$
dans le simplexe $\Delta_n = \{x \in \R^n : x_i \geq 0, \sum x_i = 1\}$.
Interpr\'eter ce r\'esultat comme un th\'eor\`eme de point fixe.
\end{exercice}
